
\needspace{3\baselineskip}
\color{uniblue}{\subsection{Токовый контроль для защиты от дуговых замыканий (ТК ЗДЗ)}}
\color{black}

\begin{enumerate}

\item
Токовый контроль для защиты от дуговых замыканий служит для формирования сигнала наличия тока через трансформатор для алгоритмов дуговых защит секций НН.

\item
В устройстве токовый контроль для ЗДЗ представлен функцией «ТК ЗДЗ».

\item
Ввод токового контроля для ЗДЗ в работу на этапе параметрирования устройства осуществляется программным переключателем SGF1 «Ввод\_функции». Информация о введенной функции отображается выходным сигналом «ТК ЗДЗ: Ввод».

\item
Функция может быть выведена из работы оперативно путем активации сигнала «ОВ ТК ЗДЗ», что характеризуется активным состоянием сигнала «ТК ЗДЗ: Оперативный вывод» на выходе функции.

\item
Функция контролирует повышение любого фазного тока ($I_A, I_B, I_C$) стороны ВН трансформатора выше заданной уставки, при превышении заданной уставки выдается пусковой сигнал для выдачи в цепи дуговых защит секций НН трансформатора.

\item
Алгоритм работы функции выполнен в соответствии с рисунком \ref{tkzdz:img1}. В таблице \ref{tkzdz:tbl1} приведены параметры, необходимые для настройки функции.

\vspace{1mm}
\begin{figure}[H]
\centering
\includegraphics[width=1\textwidth,height=1\textheight,keepaspectratio]{\fbpath/12. ТК ЗДЗ 35/_latex/img/1.pdf}
\captionof{figure}{Функциональная схема логики ТК ЗДЗ}\label{tkzdz:img1}
\end{figure}

\FloatBarrier % Все float'ы ДО этой строки будут размещены здесь.
\end{enumerate}

